\section{A single uniform bounded algorithm}
\label{sec:implementation}

We now implement the ideal choices, bound every stored structure on
every execution, and finish the proof of Theorem~\ref{thm:main}.
The program runs both trials; the deficiency split selects a success
proof, not a branch that the program must recognize. Its sampling
limits and operation counters enforce finite running time even when
the resource event of Section~\ref{sec:resources} fails.

The full schedule is as follows. On either exceptional case in the
guard, enumerate all indexed subsets exactly. Otherwise perform
\(u^{120}\) rounds. In each round run the compressed-block trial and
then the representation trial, sequentially, using fresh randomness
and reset trial state. Abort each trial after
\(u^{500}2^{\lceil n/2\rceil}\) underlying word operations; a sampling
failure has the same effect as an abort. A YES from either trial stops
the program. If no trial returns YES, return NO after the last round.

\subsection{Bounded sampling and its exact law}

An ideal draw is uniform on its specified finite set, using fresh
randomness conditional on the preceding choices. The implementation
uses a fixed attempt limit. We couple it to the ideal experiment,
rather than condition the whole experiment on every sampler succeeding:
that latter conditioning could bias earlier adaptive choices.

\begin{lemma}[Bounded samplers]\label{lem:samplers}
A uniform integer in \(\{0,\ldots,D-1\}\), with \(1\le D<2^w\),
can be sampled using at most \(n^3\) rejection attempts, or the trial
can report failure. Conditional on any preceding history and on this
call accepting, its result is uniform in the specified range. It
couples to an ideal draw with failure probability at most
\(2^{-n^3}\).

For \(4\le\ell\le n\), the same attempt limit gives a uniform prime
in \(\mathcal P_\ell\), conditional on this call accepting, and
coupling failure at most \(e^{-n^2/4}\). A full trial, including
permutations, signs, residues and the four separating families, couples
to the ideal experiment with total failure
\(\exp(-\Omega(n^2))\). All constants are absolute. All permitted
attempts are included in the deterministic work bounds.
\end{lemma}

\begin{proof}
For \(D=1\), return zero without randomness. Otherwise take
\(\lceil\log D\rceil\) low bits of a fresh random word and reject
values at least \(D\). Each value in the range has the same acceptance
probability, and an attempt succeeds with probability at least \(1/2\).
The first accepted value is therefore uniform, including conditional
on accepting within \(n^3\) attempts. Comparing the same random
attempts with an ideal repeat-until-success sampler gives failure
probability at most \(2^{-n^3}\). Computing the bit length and each
attempt takes polynomially many charged word operations. Uniform
permutations use successive range draws among the remaining positions;
a uniform fixed-size subset takes the required initial part of such a
permutation. Sign masks use fresh unbiased bits.

For a prime, draw \(\ell\) low bits, reject integers below two or
outside \(4\ell p\ge2^\ell\), and test primality by deterministic
trial division. Every prime in \(\mathcal P_\ell\) has the same
probability of being drawn. Lemma~\ref{lem:prime} gives acceptance
probability at least \(1/(4\ell)\), so the attempt limit fails with
probability at most
\[
 (1-1/(4\ell))^{n^3}\le e^{-n^3/(4\ell)}\le e^{-n^2/4}.
\]
The result is uniform conditional on this call accepting. In every
application \(\ell\le.258n\), so trial division costs at most
\(\poly(n)2^{.129n}\) operations per attempt, and its \(n^3\)
attempts preserve that form of bound.

The four families are the only exponentially numerous sampled objects.
Each family has at most \(10\cdot2^g+1\) sets, each generated by at
most \(g\) range calls. All other permutations, guesses and residues
use polynomially many range calls, and there are constantly many prime
calls. Thus each trial uses at most \(\poly(n)2^g\) range calls.
The failure bounds hold conditional on every history, including
adaptively chosen ranges. Summing them over the bounded number of
calls gives
\[
 \poly(n)2^g2^{-n^3}+O(1)e^{-n^2/4}
   =\exp(-\Omega(n^2)).
\]
This bounds the event that the bounded and ideal executions first
diverge at an exhausted sampler. Retaining only current sampler state
uses no uncounted random tape. The bounded cost, including all rejected
attempts, is the cost charged by the trial analyses.
\end{proof}

On the main branch, \(\log u\le n/10^9\), so the coupling loss is
smaller than any fixed multiple of \(u^{-110}\) for large \(n\).
A watchdog interruption does not create an additional probabilistic
assumption. On successful coupling, all sampled choices agree with the
ideal choices, while the cost estimates already allow the worst
permitted number of attempts. The favorable work events therefore
complete before the watchdog; other executions may simply abort.

\subsection{Exact word operations}

The program scans the input and counts magnitude bits by shifts to
obtain \(b\). Actual signed weights, the transformed target and all
relevant differences have magnitude less than
\[
 2^{b+\lceil\log(n+2)\rceil+4}.
\]
Even an invalid parent tuple uses an input index at most twice, so the
same bound covers the quantities considered before disjointness is
checked. Store a sign and a magnitude, with canonical zero. Signed
addition, comparison and normalized remainders are implemented on
these fields. When combining an actual weight with a residue, reduce
the weight first. A fixed-width signed weight followed by the relevant
mask bits supplies the trie key.

Every mask has at most \(n\) bits, and every stored item has a constant
number of word fields. Binomial coefficients on at most \(n\) elements
can be computed multiplicatively; immediately before division the
intermediate integer is at most \(n2^n\). Floor logarithms are bit
lengths. Hence the background allocations, modulus parameters and
rational floors are exact integer computations. The separator quantity
\(3^m\) fits in a word as well. Feasibility tests precede unsigned
subtractions, or those quantities are represented with a sign.
No entropy or real-arithmetic oracle is used by the program.

The largest product modulus satisfies \(P\le2^{s+.258n}\). Residue
arithmetic therefore fits in a word, and trial division compares the
candidate divisor with the corresponding integer quotient instead of
forming a potentially unnecessary square. Masks for range sampling
use fewer than \(w\) bits. The main-branch guard also gives
\[
 \log\bigl(u^{500}2^{\lceil n/2\rceil}\bigr)
 \le (1/2+500/10^9)n+1<w,
\]
for large \(n\). Thus fuel, the \(u^{120}\) round count, intermediate
powers, timestamps and counters have bounded word width. The compact
memory addresses fit by the space calculation below. All arithmetic
invoked as a unit-cost operation stays within the stated word model.

Permutations, sign masks, moduli, residues, family masks and replay
state remain in the counted workspace whenever reused. Once a random
value is discarded, the program never reads it again.

\subsection{Space on every random outcome}

We specify the simultaneous storage terms, rather than infer space
from the favorable resource event. Let \(L_{\max}=2^{.185n}\), an
upper bound on every raw representation leaf list on the main branch.
There are only eight such lists; an item, heap cell or stream-row
state occupies constantly many words. Linear auxiliary arrays suffice
for sorting. Hence the stored leaves, heaps, row states and a constant
number of replay snapshots occupy \(O(L_{\max})\) words. The tie
buffer uses \(O(B)\) words. Initial compressed-block lists and
mixer-counting half lists have sizes \(2^s\le2^{.12n}\); the
compressed trial's stored sublists have size at most \(2^{.14n}\).
These fit the same bound, and the two trials run sequentially.

A dictionary keeps at most \(C\) distinct records and stops when a
\((C+1)\)st key is encountered. Its fixed-width keys have \(O(u)\)
bits, so the binary trie, including transient nodes for that rejected
key, has \(O(u(C+1))\) nodes. Each node has constantly many word
fields. At most a constant number of dictionaries and retained parent
arrays coexist. In the right prefix search, an overflowing task's
workspace is discarded before its children are processed. An
\(O(n)\)-depth stack stores prefixes and return information, not all
previous task dictionaries. The completed left array is retained
while the current right task or final check runs.

For the final check, the four family-index trees have total active
depth \(O(n)\); every depth keeps constantly many arrays of length at
most \(C\). These are actual \(O(nC)\) retained record buffers,
including linear sorting workspace, as in \eqref{sep:workspace}.
Caller-retained parent arrays contribute another \(O(C)\), already
covered by this bound. The incidence tables have eight lists per
block mask across the four blocks, each of length at most
\(10\cdot2^g+1\). Their total capacity is at most
\[
 8\,2^g(10\cdot2^g+1)=O(2^{2g}).
\]
Each stored index uses at most \(g+4\) bits, and each family mask uses
\(g\) bits, so both fit in one word. Directories and sampled families
add only \(O(2^g)\) words. The incidence contribution has the smaller
exponential factor
\(2^{2g}\le2^{.12n}\).

Combination generation, permutations, sampler state, integer
parameters, prefix and interval control, counters, and a conservative
binomial workspace use \(O(u^2)\) further words. In particular,
none retains the hypothetical check list. Combining all terms, some
absolute implementation constant \(K\) bounds simultaneous storage by
\[
 K\bigl[u(C+1)+(n+1)C+L_{\max}+B+2^{2g}+u^2\bigr].
\]
Now \(C\le2^{.199n}\), \(L_{\max}\le2^{.199n}\),
\(B\le2^{.199n}\), \(2^{2g}\le2^{.199n}\), and \(n\le u\).
Term by term this proves the required conservative bound
\begin{equation}\label{eq:space-final}
 O\bigl(u^{15}2^{.199n}\bigr)
\end{equation}
on \emph{every} random outcome, including an overflowing or aborted
trial. No assumption about average bucket size or \(\mathcal E\) was
used. Its address bit length is at most
\(.199n+15\log u+O(1)<w\) on the main branch.
Allocation and stack pointers are reset between tasks and trials;
entries are initialized when reused. Even clearing the allocated
storage linearly is below the trial time budget, so aborts accumulate
neither stored data nor uncharged cleanup work.

\subsection{A finite program, amplification and the theorem}

There are finitely many sufficiently-large conditions in the argument,
with absolute constants and fixed polynomial degrees. The main-branch
guard makes each fixed exponential margin dominate the required
polynomial factors uniformly in that branch. Choose one absolute
integer \(n_0\) meeting all these conditions and fix it in the program,
independently of \(c\).

Small-word inputs need not hold \(n_0\) as a word. Place an exact
subset-enumeration routine near the beginning of the finite program.
Unroll the binary digits of \(n_0\), from least to most significant,
using only literals zero and one. Maintain a shifted copy of \(n\)
and a three-valued comparison state. At each digit, compare its low
bit with the fixed threshold digit, replace the comparison state
when they differ, and shift the copy right. A newly seen difference
has higher significance. If the copy becomes zero before the
threshold's leading digit has been processed, branch to the nearby
exhaustive routine. After that leading digit, a nonzero remaining copy
means \(n>n_0\); otherwise the comparison state decides.

A short input executes only \(O(1+\log n)\) adjacent digit stages
before exiting. Its instruction addresses and state therefore remain
small; it never loads the full threshold or a distant main-program
address. Once \(n\ge n_0\), enlarge this fixed threshold if necessary
so that all other fixed literals and main-program addresses fit.
The second guard, \eqref{eq:guard}, is then an integer comparison.
Both exceptional cases use exact enumeration with an \(n\)-bit subset
counter and \(O(n)\) words. This is one finite program with no advice
and no input for the exponent \(c\).

On a main-branch YES input, Lemma~\ref{lem:compressed} gives constant
success probability when \(\mu\ge.00765n\), and
Lemma~\ref{lem:resources} gives probability at least \(u^{-110}\)
when \(\mu<.00765n\). Both trials run in each round with fresh words
and reset state. Thus each round succeeds with probability at least
\(u^{-110}\), also conditional on every previous round failing.
After \(u^{120}\) rounds the probability of missing a solution is at
most
\[
 (1-u^{-110})^{u^{120}}\le e^{-u^{10}}<1/3.
\]
Every YES is an exact valid match by Lemma~\ref{alg:validity}; the
sign-toggle identity certifies an original solution. The exceptional
branch is exact, so correctness holds on every individual input.

Instrument each trial with a counter for underlying word operations.
Its constant overhead, stopping and reset give per-trial time
\(O(u^{500}2^{\lceil n/2\rceil})\) on every outcome. Input scanning
and the initial guards take polynomial time in \(u\). All trial
preprocessing, sorting, table construction and sampling are charged.
The total main-branch time is consequently
\[
 O\bigl(u^{620}2^{\lceil n/2\rceil}+\poly(u)\bigr).
\]
For every fixed positive integer \(c\), all sufficiently large inputs
with \(b\le n^c\) avoid the exceptional branch and have
\(u=\poly_c(n)\). This time bound and \eqref{eq:space-final} hold
simultaneously for every random outcome. Together with the correctness
argument they prove Theorem~\ref{thm:main}. Finally, for each fixed
\(\rho>.199\), the factor \(u^{15}=\poly_c(n)\) is absorbed by
\(2^{(\rho-.199)n}\), proving Corollary~\ref{cor:target}; the
polynomial factor remains at the endpoint \(.199\).
